%%%
%%% DO NOT EDIT - Applicants do not (usually) need to touch this file
%%%

%%%
%%% Application form LaTeX template for the computing facilities at CfCA/NAOJ
%%%
%%% 2024/05/31 Takashi Ito (revised for the 2024 fall application for XD2000)
%%% 2021/12/14 Takashi Ito (revised for the 2022 application)
%%% 2021/06/15 Takashi Ito (substantially revised again)
%%% 2020/06/10 Takashi Ito (substantially revised)
%%% 2020/02/28 Kaori Kano  (slightly revised)
%%% 2020/02/08 Takashi Ito (initial draft)
%%%

\documentclass[11pt,a4j]{jarticle}

%%% Abbreviations for journals - recommendation by Tomoya Takiwaki
\usepackage{aas_macros}

%\pagestyle{plain}

% Format version of the template files
%\def\formversion{2021--04--01--cfca}
%\def\formversion{2021--10--01--cfca}
%\def\formversion{2022--04--01--cfca}
%\def\formversion{2022--04--01--cfca (2)}
 \def\formversion{2024--07--01--cfca}

% Definition of the colors for the remarks
\def\remarkcolor{blue}
\def\noneedcolor{red}

% Do not touch
\setcounter{secnumdepth}{2}
\setcounter{tocdepth}{2}
\setlength{\textheight}{250mm}
\setlength{\topmargin}{-15mm}

% In future this may be changed...
%\def\CfCASUPERCOMP{XC50}
 \def\CfCASUPERCOMP{XD2000}

% Vertical spacing
\def\vspaceplussmall{\vspace{0.3\baselineskip}}
\def\vspaceminussmall{\vspace{-0.3\baselineskip}}
\def\vspaceforexamples{\vspace{0.36\baselineskip}}

% For you to include EPS figures...
\usepackage{graphicx}
%\usepackage[dvipdfmx]{graphicx} % If you need, use the option

% We always need some fun...
\usepackage{color}

% Hope most users' environment has them...
\usepackage{url}

%%%
%%% When you include DOI entries in your *.bib file, you may need to uncommenet the following lines, provided your environment is ready - do it at your own risk.
%%%
%\usepackage{xcolor}
%\usepackage[dvipdfmx]{hyperref}
%\hypersetup{
%    colorlinks,
%    linkcolor={red!50!black},
%    citecolor={blue!50!black},
%    urlcolor={blue!80!black}
%}

% See for detail: https://www.okomeda.net/wp/434/
\usepackage{fancyhdr} 
\pagestyle{fancy} 
%\lhead[pageeven]{pageodd}
%\chead[pageeven]{pageodd}
\rhead[]{\color[gray]{0.5}{\footnotesize {\formversion}}}
%\lfoot[pageeven]{pageodd}
\cfoot[]{\thepage}
%\rfoot[pageeven]{pageodd}
\renewcommand{\headrulewidth}{0pt}
\renewcommand{\footrulewidth}{0pt}

% Comment multiple lines
\long\def\comment#1{}

% Place the table title at the vertical center of a column by lowering it
\newcommand{\lw}[1]{\smash{\lower2.ex\hbox{#1}}}
%
% Place the table title at the vertical center of a column by lowering it more
\newcommand{\llw}[1]{\smash{\lower4.ex\hbox{#1}}}

%%%
%%% Flag for the English version
%%%
\newif\ifEnglishform
%\Englishformtrue
 \Englishformfalse

%%%
%%% Definition of the if structures
%%%
\newif\iffinalform
%\finalformtrue
 \finalformfalse

%%%
%%% When you use BibTeX
%%%
\newif\ifUseBibTeX
%\UseBibTeXtrue
 \UseBibTeXfalse

%%%
%%% Definition of categories
%%%
\newif\ifXDS
%\XDStrue
 \XDSfalse

\newif\ifXDA
%\XDAtrue
 \XDAfalse

\newif\ifXDB
%\XDBtrue
 \XDBfalse

\newif\ifXDMD
%\XDMDtrue
 \XDMDfalse

\newif\ifXDTrial
%\XDTrialtrue
 \XDTrialfalse

\newif\ifGPPC
%\GPPCtrue
 \GPPCfalse

\newif\ifGPU
%\GPUtrue
 \GPUfalse

\newif\ifSPC %iwasaki
%\SPCtrue
 \SPCfalse

% User's choice
\input{parameters}

% Which category this document is for
\ifXDS
  \def\カテゴリ{XD--S}
\fi
\ifXDA
  \def\カテゴリ{XD--A}
\fi
\ifXDB
  \def\カテゴリ{XD--B}
\fi
\ifXDMD
  \def\カテゴリ{XD--MD}
\fi
\ifXDTrial
  \def\カテゴリ{XD--Trial}
\fi
\ifGPPC
  \ifEnglishform
    \def\カテゴリ{GP--PC (general-purpose PC farm)}
  \else
    \def\カテゴリ{GP--PC (計算サーバ)}
  \fi
\fi
\ifGPU
  \ifEnglishform
    \def\カテゴリ{GPU cluster}
  \else
    \def\カテゴリ{GPUクラスタ}
  \fi
\fi
\ifSPC %iwasaki
  \ifEnglishform
    \def\カテゴリ{Small Parallel Computers}
  \else
    \def\カテゴリ{中規模サーバ}
  \fi
\fi

\ifEnglishform
  \def\nopageexeed{(The amount of description should not exceed one page)}
\else
  \def\nopageexeed{(記載はこのページの範囲を超えないようにしてください)}
\fi

%
% Example resource values for the XD-* categories
%
\iffinalform
  \def\Mnodetime{--}   % node*time [k nodes hours] for M
  \def\Mmempernode{--} % memory/node [GB/node]     for M
  \def\Pnodetime{--}   % node*time [k nodes hours] for P
  \def\Pmempernode{--} % memory/node [GB/node]     for P
\else
  \def\Mnodetime{8.4}  % node*time [k nodes hours] for M
  \def\Mmempernode{50} % memory/node [GB/node]     for M
  \def\Pnodetime{33.6} % node*time [k nodes hours] for P
  \def\Pmempernode{200}% memory/node [GB/node]     for P
\fi

%%%
%%% When you use BibTeX for citing previous literature
%%%
\ifUseBibTeX
  \usepackage[numbers]{natbib} % or whatever you like
 %\usepackage{natbib}          % or whatever you like
  \setlength{\bibsep}{-1.6pt}
  % Changing the figure variable and \refname in the English form
  \ifEnglishform
    \makeatletter
    \def\fnum@figure{Figure\ \thefigure}
    \makeatother
    \renewcommand{\refname}{\normalsize References cited in Sections \ref{sec:objective}--\ref{sec:codeprep}}
    \renewcommand{\refname}{}  % Remove the ref section title
  \else
    %%% Not necessary to redefine it in the Japanese version
    %\renewcommand{\refname}{\normalsize 第\ref{sec:objective}--\ref{sec:codeprep}節で引用された文献}
    \renewcommand{\refname}{}  % Remove the ref section title
  \fi
\else
  % Do not use BibTeX
\fi

%%% End of the BibTeX part

%%%%%%%%%%%%%%%%%%%%%%%%%%%%%%%%%%%%%%%%%%%%%%%%%%%%%%%%%%%%%%%%%%%%%%%
\begin{document}

\vspace*{-2\baselineskip}
\begin{center}
\ifEnglishform
   {\huge {\textbf{{\カテゴリ} application form}}}
\else
   {\huge {\textbf{{\カテゴリ} 利用申請書}}}
\fi
\end{center}

%%%
%%% 研究目的
%%%
\iffinalform
% Do nothing
\else
  \ifXDS
    \ifEnglishform
      \textcolor{\remarkcolor}{This application form is for the large-scale execution category of {\CfCASUPERCOMP} (XD--S). You cannot use it for applications in normal categories (XD--A, XD--B, XD--MD, XD--Trial).}
    \else
      \textcolor{\remarkcolor}{本申請書は{\CfCASUPERCOMP}の大規模実行カテゴリ(XD--S)へ申請するためのものです。通常のカテゴリ(XD--A, XD--B, XD--MD, XD--Trial)への申請には使えませんので、ご注意ください。}
    \fi % End of \ifEnglishform
  \else
    \ifGPPC
      \ifEnglishform
        \textcolor{\remarkcolor}{This application form is exclusively for {\カテゴリ}. You cannot use it to apply for the other categories such as {\CfCASUPERCOMP} or GPU.}
      \else
        \textcolor{\remarkcolor}{本申請書は{\カテゴリ}へ申請するためのものです。{\CfCASUPERCOMP}やGPUなど他カテゴリへの申請には使えませんので、ご注意ください。}
      \fi % End of \ifEnglishform
    \else
      \ifGPU
        \ifEnglishform
          \textcolor{\remarkcolor}{This application form is exclusively for {\カテゴリ}. You cannot use it to apply for the other categories such as {\CfCASUPERCOMP} or GP--PC (general-purpose PC farm).}
        \else
          \textcolor{\remarkcolor}{本申請書は{\カテゴリ}へ申請するためのものです。{\CfCASUPERCOMP}やGP--PC (計算サーバ)など他カテゴリへの申請には使えませんので、ご注意ください。}
        \fi % End of \ifEnglishform
      \else
        \ifSPC
          \ifEnglishform
              \textcolor{\remarkcolor}{This application form is exclusively for {\カテゴリ}. You cannot use it to apply for the other categories such as {\CfCASUPERCOMP} or GP--PC (general-purpose PC farm).}
          \else
            \textcolor{\remarkcolor}{本申請書は{\カテゴリ}へ申請するためのものです。{\CfCASUPERCOMP}やGP--PC (計算サーバ)など他カテゴリへの申請には使えませんので、ご注意ください。}
          \fi % End of \ifEnglishform
        \else
          \ifEnglishform
            \textcolor{\remarkcolor}{This application form is for the regular categories of {\CfCASUPERCOMP} (XD--A, XD--B, XD--MD, and XD--Trial). You cannot use it to apply for the category XD--S.}
          \else
            \textcolor{\remarkcolor}{本申請書は{\CfCASUPERCOMP}の通常のカテゴリ(XD--A, XD--B, XD--MD, XD--Trial)へ申請するためのものです。カテゴリXD--Sへの申請には使えませんので、ご注意ください。}
          \fi % End of \ifEnglishform
        \fi % End if \ifSPC
      \fi % End of \ifGPU
    \fi % End of \ifGPPC
  \fi % End of \ifXDS
\fi % End of \iffinalform

%%%%%%%%%%%%%%%%%%%%%%%%%%%%%%%%%%%%%%%%%%%%%%%%%%%%%%%%%%%%%%%%%%%%%%%
\ifEnglishform
  \section{Significance and objective of your research\label{sec:objective}}
\else
  \section{研究の意義と目的\label{sec:objective}}
\fi

\iffinalform
  \input{1-objective}
\else
  \ifEnglishform
    \textcolor{\remarkcolor}{Describe the significance and objective of your research. Clarify the background, scientific goal, and originality of your approach.
}
    \textcolor{\remarkcolor}{\\ \hspace*{1em}
Provide the bibliographic information for the references cited in Section \ref{sec:objective}, \ref{sec:method}, \ref{sec:codeprep} in Section \ref{sec:citedreference}.
The format of the bibliography and how to cite them in the text should follow the AAS Journals.
Examples of in-text citations:
``There are several TNOs with perihelion distance greater than 45 au that have clearly experienced planetary scattering \citep{bannister2017}, and we need to examine the results of a larger range of observations to verify them \citep[e.g.][]{boe2019,grav2011}.''
The above examples show citations by reference number, but citations in author$+$year format such as ``(\citeauthor{grav2011} \citeyear{grav2011}; \citeauthor{boe2019} \citeyear{boe2019})'' are also acceptable.
}
\textcolor{\remarkcolor}{\\ \hspace*{1em}
The {\LaTeX} version of the application form requires \textsc{Bib}{\TeX} in order to create the list of references
(note that you will need to prepare the necessary environment to run \textsc{Bib}{\TeX}).
Follow the instructions in the \texttt{00READMEe.txt} file.
Make sure that the amount of text does not exceed the page limit for each section.
}
% \textcolor{\remarkcolor}{\\ \hspace*{1em}
% When you cite references, include (embed) the bibliographic information in the text. This is all the same in the following sections.
% Example of citations: ``There are several TNOs with perihelion distance greater than 45 au that have clearly experienced planetary scattering (Bannister et al. 2017, AJ, 153, 262).''
% You may also use the number reference system (i.e. inline citation is numbered as [1] or [2,~4]), and all reference sources are listed somewhere in p. \pageref{sec:objective}--\pageref{sec:codeprep}.
% The {\LaTeX} version of the application file allows you to use \textsc{Bib}{\TeX} for this purpose.
% Follow the instructions in the file \texttt{00READMEe.txt}
% (but you must prepare your own environment for \textsc{Bib}{\TeX}).
% Be sure that the amount of information in each section does not exceed the page limit.
% }
  \else
    \textcolor{\remarkcolor}{研究の意義や目的を具体的かつ分かりやすくまとめてください。その際には研究の学術的な背景、今回の研究で明らかにしようとしている点、特色（独自性）が分かるように記載してください。}
    \ifGPPC
    % Do nothing for GP-PC
    \else
      \ifGPU
      % Do nothing for GPU
      \else % Otherwise (XD2000), do things
        \textcolor{\remarkcolor}{また、CfCAホームページ上の{\CfCASUPERCOMP}募集要項 \url{https://www.cfca.nao.ac.jp/xd2000application} にある「7. 審査」の内容も参考にしてください。}
      \fi % End of \ifGPU
    \fi % End if \ifGPPC
    \textcolor{\remarkcolor}{\\ \hspace*{1em}
第\ref{sec:objective}, \ref{sec:method}, \ref{sec:codeprep}節で引用した文献の書誌情報を第\ref{sec:citedreference}節に記してください。
文中での引用方法と文献一覧 (書誌情報) の形式はAAS Journalsに準じます。
文中での引用の例：
「近点距離が45 auを超えるTNOには明らかに惑星散乱を経験した天体が幾つも含まれており \citep{bannister2017}、
  それらの実証のためにはより広範囲にわたる観測結果を精査する必要がある \citep[例.][]{boe2019,grav2011}。」
上記の例は文献番号による引用ですが、
``\citeauthor{bannister2017} (\citeyear{bannister2017})''や
``(\citeauthor{grav2011} \citeyear{grav2011}; \citeauthor{boe2019}  \citeyear{boe2019})'' といった著者名$+$刊行年形式の引用でも構いません。
}
\textcolor{\remarkcolor}{\\ \hspace*{1em}
{\LaTeX}版の申請書では文献引用のために\textsc{Bib}{\TeX}が必要です (但し\textsc{Bib}{\TeX}の動作環境は各人でご準備ください)。
\texttt{00READMEj.txt}ファイル内の指示に従って記載を行ってください。
この場合も、各節の記載量がそれぞれのページ数制限を超えないよう留意してください。
}
  \fi % End if \ifEnglishform
  \vfill
  \begin{flushright}
  \textcolor{\remarkcolor}{\small \nopageexeed}
  \end{flushright}
\fi % End of %iffinalform

%%%
%%% Dealing with page number excess
%%%
\def\npagemax{1}
\def\seclabelname{\ref{sec:objective}}
% Common to all the sections
\ifnum\value{page}>\npagemax%
  \vspace{1.0\baselineskip}
  \noindent
  \ifEnglishform
    \textcolor{red}{\textbf{\Large \fbox{Error!} Section {\seclabelname} exceeds the page limit (up to p.~{\npagemax}).}}
    \typeout{===ERROR===}
    \typeout{===ERROR=== Section \seclabelname exceeds the page limit (up to p.\npagemax).}
    \typeout{===ERROR===}
  \else
    \textcolor{red}{\textbf{\Large \fbox{警告} 第{\seclabelname}節がページ数の上限(p.~{\npagemax}まで)を超えています。(\thepage)}}
    \typeout{===警告===}
    \typeout{===警告=== 第\seclabelname 節がページ数の上限(p.\npagemax まで)を超えています。}
    \typeout{===警告===}
  \fi
\else
  % No problem in the page number limitation...
\fi

\clearpage
%%%%%%%%%%%%%%%%%%%%%%%%%%%%%%%%%%%%%%%%%%%%%%%%%%%%%%%%%%%%%%%%%%%%%%%
\ifEnglishform
  \section{Research plan and methods\label{sec:method}}
\else
  \section{研究計画と方法\label{sec:method}}
\fi

\iffinalform
  \input{2-method}
\else
  \ifEnglishform
    \textcolor{\remarkcolor}{Summarize your research plan and methods here. Describe the specific outcomes that you expect to achieve during this fiscal year (ending in March). If there are any new or unique perspectives on your research methods, specify them as well.
The format of the citation must follow the instruction specified in Section \ref{sec:objective}.
}
  \else
    \textcolor{\remarkcolor}{研究計画と方法を具体的に記してください。その際には本年度中の達成が期待される具体的な成果を分かりやすく述べてください。もし研究方法に新しい視点や独自性があれば、それも明記してください。
文献引用の形式は第\ref{sec:objective}節の説明に記した通りです。}
  \fi % End of \ifEnglishform
  \vfill
  \begin{flushright}
  \textcolor{\remarkcolor}{\small \nopageexeed}
  \end{flushright}
\fi % End of \iffinalform

%%%
%%% Dealing with page number excess
%%%
\def\npagemax{2}
\def\seclabelname{\ref{sec:method}}
% Common to all the sections
\ifnum\value{page}>\npagemax%
  \vspace{1.0\baselineskip}
  \noindent
  \ifEnglishform
    \textcolor{red}{\textbf{\Large \fbox{Error!} Section {\seclabelname} exceeds the page limit (up to p.~{\npagemax}).}}
    \typeout{===ERROR===}
    \typeout{===ERROR=== Section \seclabelname exceeds the page limit (up to p.\npagemax).}
    \typeout{===ERROR===}
  \else
    \textcolor{red}{\textbf{\Large \fbox{警告} 第{\seclabelname}節がページ数の上限(p.~{\npagemax}まで)を超えています。(\thepage)}}
    \typeout{===警告===}
    \typeout{===警告=== 第\seclabelname 節がページ数の上限(p.\npagemax まで)を超えています。}
    \typeout{===警告===}
  \fi
\else
  % No problem in the page number limitation...
\fi

\clearpage

%%%%%%%%%%%%%%%%%%%%%%%%%%%%%%%%%%%%%%%%%%%%%%%%%%%%%%%%%%%%%%%%%%%%%%%
\ifEnglishform
  \section{Preparation of computational code\label{sec:codeprep}}
\else
  \section{計算コードの準備状況\label{sec:codeprep}}
\fi

\iffinalform
% Do nothing
\else
  \ifEnglishform
    \textcolor{\remarkcolor}{({\S}\ref{ssec:codedevel} and {\S}\ref{ssec:requiredres} combined, do not exceed one page.)}
  \else
    \textcolor{\remarkcolor}{({\S}\ref{ssec:codedevel}と{\S}\ref{ssec:requiredres}を合わせて1ページ以内)}
  \fi
\fi

%%%
%%% 3.1
%%%
\ifEnglishform
  \subsection{Development and optimization of computational codes\label{ssec:codedevel}}
\else
  \subsection{計算コードの開発と最適化状況\label{ssec:codedevel}}
\fi

\def\myfirstcaseO{%
\ifEnglishform
  \textcolor{\noneedcolor}{\fbox{Applications in {\カテゴリ} do not require any descriptions of this section.}}
\else
  \textcolor{\noneedcolor}{\fbox{このカテゴリ({\カテゴリ})への申請では本節へ記載は不要です}}
\fi
} % End of \def\myfirstcaseO{%

\def\mysecondcaseO{%
\ifEnglishform
  \textcolor{\remarkcolor}{Describe the current status of the development and optimization of your numerical code.
If you are a former user of the computers at CfCA, you need to clarify that the numerical code is expected to run efficiently on this system by showing the scalability of the code.
Particularly when you apply for XD--A,
%or XD--S,
you need to show the scalability of at least up to 1,000 cores.
The format of the citation must follow the instruction specified in Section \ref{sec:objective}.
}
\else
  \textcolor{\remarkcolor}{計算コードの開発と最適化の状況を具体的に記してください。
過去にCfCAの共同利用計算機を利用した経験があればその状況も含めてください。
スケーラビリティ等の実測値を用い、効率の良い計算が実施可能であることを示してください。
特にXD--A
%またはXD--S
に申請する場合は、1000並列以上を用いた場合のスケーラビリティを示してください。
文献引用の形式は第\ref{sec:objective}節の説明に記した通りです。}
\fi
} % End of \def\mysecondcaseO{%

\ifXDS
  \iffinalform
    \input{3-1-codedevel}
  \else
    \mysecondcaseO
  \fi
\fi
\ifXDA
  \iffinalform
    \input{3-1-codedevel}
  \else
    \mysecondcaseO
  \fi
\fi
\ifXDB
  \iffinalform
    \input{3-1-codedevel}
  \else
    \mysecondcaseO
  \fi
\fi
\ifXDMD
  \iffinalform
    \input{3-1-codedevel}
  \else
    \mysecondcaseO
  \fi
\fi
\ifXDTrial
    \myfirstcaseO
\fi
\ifGPPC
    \myfirstcaseO
\fi
\ifGPU
    \myfirstcaseO
\fi
\ifSPC
    \myfirstcaseO
\fi

%%%
%%% 3.2
%%%
\ifEnglishform
  \subsection{Estimate of the amount of required computing resources\label{ssec:requiredres}}
\else
  \subsection{必要な計算資源の見込み\label{ssec:requiredres}}
\fi

\def\myfirstcaseO{%
\ifEnglishform
  \textcolor{\noneedcolor}{\fbox{Applications in {\カテゴリ} do not require any descriptions of this section.}}
\else
  \textcolor{\noneedcolor}{\fbox{このカテゴリ({\カテゴリ})への申請では本節へ記載は不要です}}
\fi
} % End of \def\myfirstcaseO{%

%%%
%%% Along the request from Tomoya Takiwaki on June 18, 2021 (ref. #10066)
%%%
\def\reqresourceXD{%
%%%
%%% English
%%%
\ifEnglishform
{\small
  \textcolor{\remarkcolor}{
Describe the estimated amount of computing resources required to carry out your proposal (number of nodes, number of CPU cores, amount of memory, the time required for one model calculation, the total computation time for all the model calculations, etc.).
``k'' denotes 1,000 (ex. 200 k $= 200,000$) in what follows. 
Provide the values for each system you wish to use (System M and System P).
\begin{quote}
(Example)
It takes a week  to execute a run of our simulation using 10 nodes on System M.
With 5 model calculations performed, the total computation time will be  8.4 k nodes hours.
Memory usage per node is  50 GB/node. \\
It takes 2 weeks to execute a run of our simulation using 20 nodes on System P.
With 5 model calculations performed, the total computation time will be 33.6 k nodes hours.
Memory usage per node is 200 GB/node.
\end{quote}
\hspace*{1em}
Note that the total computation time value should be given in the form of a node time product (in k nodes hours).
For your reference, the following information is based on the operational statistics obtained so far.
}
%
\vspaceforexamples
\textcolor{\remarkcolor}{
\begin{center}
\begin{tabular}{lrr}
\hline
\lw{\small category} & {\small appropriate node time product}         & {\small relative value to the} \\
                     & {\small / (year $\cdot$ user) (k nodes hours)} & {\small whole resource of {\CfCASUPERCOMP}}            \\
\hline
XD--A    & 50--250                                        &  2--10{\%} \\
XD--B$+$ & 25--130                                        &  1--5{\%}  \\
XD--B    & 13--60                                         &  0.5--2.5{\%} \\
XD--MD   &  6--30                                         &  0.25--1.3{\%} \\
\hline
\end{tabular}
\end{center}
}% End if \textcolor{\remarkcolor}

\vspaceforexamples
\textcolor{\remarkcolor}{
Note that here you should fill in the value per node (in GB/node) as the memory usage.
The amount of memory per node that System M has is 128 GB/node and that of System P is 512 GB/node.
Also, provide the numbers you estimated here in the table below as well.
Enter 0 for systems you will not use.
% If you use the Word version of the template file, do not change the text in black letters in the table.
% Also, do not move the position of the table.
% If you use the {{LaTeX}} version of the template file, you do not need to be aware of these issues.
}% End of \textcolor{\remarkcolor}
}% End of \small
\noderesourcetable
\else % End of the English version
%%%
%%% The Japanese verion follows from here
%%%
\textcolor{\remarkcolor}{
申請課題を実行するのに必要となる計算資源の見込み(ノード数・CPUコア数・メモリ量・ひとつのモデル計算に要する時間・全モデル分の総計算時間など)を以下の書式で記入してください。
以下で``k''は1,000を意味します (例. 200 k $= 200,000$)。 
また、各々の値は利用を希望するシステムごとに記入してください (システムMおよびシステムP)。
\begin{quote}
(例)
計算1の1モデルを実行するにはシステムMの10ノードを用いて1週間を要する。
5モデルの実施による総計算時間は8.4 k nodes hoursとなり、ノード当たりのメモリ使用量は50 GB/nodeとなる。\\
計算2の1モデルを実行するにはシステムPの20ノードを用いて2週間を要する。
5モデル計算の実施による総計算時間は33.6 k nodes hoursとなり、ノード当たりのメモリ使用量は200 GB/nodeとなる。
\end{quote}
\hspace*{1em}
ここで総計算時間の数値はノード時間積 (単位k nodes hours) の形で記入することに注意してください。
参考までに、これまでの運用統計に基づく情報を以下に記します。
}
\vspaceforexamples
\textcolor{\remarkcolor}{
\begin{center}
\begin{tabular}{lrr}
\hline
\lw{\small カテゴリ} & {\small 適正なノード時間積 (通年,} & {\small {\CfCASUPERCOMP}計算資源の} \\
                     & {\small 1名あたり, k nodes hours)} & {\small 全量と比較した値}            \\
\hline
XD--A    & 50--250                                        &  2--10{\%} \\
XD--B$+$ & 25--130                                        &  1--5{\%}  \\
XD--B    & 13--60                                         &  0.5--2.5{\%} \\
XD--MD   &  6--30                                         &  0.25--1.3{\%} \\
\hline
\end{tabular}
\end{center}
}% End if \textcolor{\remarkcolor}

\vspaceforexamples
\textcolor{\remarkcolor}{
ここではメモリ使用量としてノード当たりの値(単位はGB/node)を記入することに注意してください。
システムMが搭載するノード当たりのメモリ量は128 GB/nodeであり、システムPでは512 GB/nodeです。
また、ここで見積もった数値は以下の表内にも記してください。
使用しないシステムについては0を入力します。
% なおWord版の雛型ファイルを使う場合には、表内に黒字で書かれた文字を変更しないでください。
% また、表の位置も動かさないでください。
% {\LaTeX}版を使う場合にはこうした事項を意識する必要はありません。
}% End of \textcolor{\remarkcolor}
\noderesourcetable
\fi
} % End of \def\reqresourceXD{%

\def\reqresourceOTHER{%
\ifEnglishform
  \textcolor{\remarkcolor}{%
Describe the estimated amount of computing resources required to carry out your proposal (number of nodes, number of CPU cores, amount of memory, the time required for one model calculation, the total computation time for all the model calculations, etc.).
}
\else
  \textcolor{\remarkcolor}{%
申請課題を実行するのに必要となる計算資源の見込み(ノード数・CPUコア数・メモリ量・ひとつのモデル計算に要する時間・全モデル分の総計算時間、など)を記してください。}
\fi % End of \ifEnglishform
} % End of \def\reqresourceOTHER{%

%%%
%%% Table value color which depends on the output mode
%%%
\iffinalform
  \def\tblvcolor{black}
\else
  \def\tblvcolor{blue}
\fi

%%%
%%% Table about the node resource info in Section 3.2
%%%
\def\noderesourcetable{%
\ifEnglishform
  \begin{center}
  \begin{tabular}{lrr}
  \hline
      & \multicolumn{1}{c}{System M} & \multicolumn{1}{c}{System P} \\
  \hline
  total computation time (k nodes $\times$ hours) & \textcolor{\tblvcolor}{\Mnodetime}   & \textcolor{\tblvcolor}{\Pnodetime}   \\
  memory amount per node (GB/node)                & \textcolor{\tblvcolor}{\Mmempernode} & \textcolor{\tblvcolor}{\Pmempernode} \\
  \hline
  \end{tabular}
  \end{center}
\else
  \begin{center}
  \begin{tabular}{lrr}
  \hline
      & \multicolumn{1}{c}{システムM} & \multicolumn{1}{c}{システムP} \\
  \hline
  ノード時間積 (k nodes $\times$ hours) & \textcolor{\tblvcolor}{\Mnodetime}   & \textcolor{\tblvcolor}{\Pnodetime}  \\
  ノード当たりのメモリ量 (GB/node)      & \textcolor{\tblvcolor}{\Mmempernode} & \textcolor{\tblvcolor}{\Pmempernode} \\
  \hline
  \end{tabular}
  \end{center}
\fi % End of \def\noderesourcetable
}% End of textcolor

%%%%%%%%%%%%%%%%%
%\vfill
% \par
% \begin{center}
% \begin{tabular}{lrr}
% \hline
%       & \multicolumn{1}{c}{システムM} & \multicolumn{1}{c}{システムP} \\
% \hline
% ノード時間積 (k nodes $\times$ hours) & \Mnodetime   & \Pnodetime   \\
% ノード当たりのメモリ量 (GB/node)      & \Mmempernode & \Pmempernode \\
% \hline
% \end{tabular}
% \end{center}
%%%%%%%%%%%%%%%%%
%} % End of \def\noderesourcetable{%

\iffinalform
  \input{3-2-requiredres}
  \ifXDS
    \noderesourcetable
  \fi
  \ifXDA
    \noderesourcetable
  \fi
  \ifXDB
    \noderesourcetable
  \fi
  \ifXDMD
    \noderesourcetable
  \fi
  \ifXDTrial
    \noderesourcetable
  \fi
\else
  \ifXDS
    \reqresourceXD
  \fi
  \ifXDA
    \reqresourceXD
  \fi
  \ifXDB
    \reqresourceXD
  \fi
  \ifXDMD
    \reqresourceXD
  \fi
  \ifXDTrial
    \reqresourceXD
  \fi
  \ifGPPC
    \reqresourceOTHER
  \fi
  \ifGPU
    \reqresourceOTHER
  \fi
  \ifSPC
    \reqresourceOTHER
  \fi
  \vfill
  \begin{flushright}
  \textcolor{\remarkcolor}{\small \nopageexeed}
  \end{flushright}
\fi

%%%
%%% Dealing with page number excess
%%%
\def\npagemax{3}
\def\seclabelname{\ref{sec:codeprep}}
% Common to all the sections
\ifnum\value{page}>\npagemax%
  \vspace{1.0\baselineskip}
  \noindent
  \ifEnglishform
    \textcolor{red}{\textbf{\Large \fbox{Error!} Section {\seclabelname} exceeds the page limit (up to p.~{\npagemax}).}}
    \typeout{===ERROR===}
    \typeout{===ERROR=== Section \seclabelname exceeds the page limit (up to p.\npagemax).}
    \typeout{===ERROR===}
  \else
    \textcolor{red}{\textbf{\Large \fbox{警告} 第{\seclabelname}節がページ数の上限(p.~{\npagemax}まで)を超えています。(\thepage)}}
    \typeout{===警告===}
    \typeout{===警告=== 第\seclabelname 節がページ数の上限(p.\npagemax まで)を超えています。}
    \typeout{===警告===}
  \fi
\else
  % No problem in the page number limitation...
\fi

\clearpage
%%%%%%%%%%%%%%%%%%%%%%%%%%%%%%%%%%%%%%%%%%%%%%%%%%%%%%%%%%%%%%%%%%%%%%%
\ifEnglishform
 %\section{Cited referenes\label{sec:citedreference}}
  \section{References cited in Sections \ref{sec:objective}--\ref{sec:codeprep}\label{sec:citedreference}}
\else
  \section{第\ref{sec:objective}--\ref{sec:codeprep}節で引用された文献\label{sec:citedreference}}
\fi

\iffinalform
% \vspace{-1.5\baselineskip}
  \bibliographystyle{aasjournal}
  \input{4-citedrefs}
\else
  \ifEnglishform
\textcolor{\remarkcolor}{%
List the references cited in Sections \ref{sec:objective}--\ref{sec:codeprep}.
Your own research achievement must be listed in Section \ref{sec:pastwork}.
However, if any of them are cited in Sections \ref{sec:objective}--\ref{sec:codeprep}, they should also be listed in this section.
Bibliographic information must be complete, and at least include the author name, journal (book) name, volume number (if journal), page number (or article id), and publication year.
You may want to add the paper title and DOI of the publication.
}
  \else
\textcolor{\remarkcolor}{%
ここには第\ref{sec:objective}節から第\ref{sec:codeprep}節の記載に於いて引用した文献を記してください。
申請者自身の従来の研究成果は第\ref{sec:pastwork}節に記入して頂きますが、それらのうち第\ref{sec:objective}--\ref{sec:codeprep}節で引用されているものは本節にも記す必要があります。
文献の書誌情報には少なくとも以下を含んでください： 著者名、雑誌名 (または書籍名)、巻号 (雑誌の場合)、ページ番号、出版年。
論文のタイトルやDOIを含んでも構いません。
}
  \fi
%%% Common to EN and JP in the example description
  \vspace{-3.0\baselineskip}     % Do you want to change it?
  \bibliographystyle{aasjournal}
  \textcolor{\remarkcolor}{
  \bibliography{bibsamples}%
  }

  \vfill
  \begin{flushright}
  \textcolor{\remarkcolor}{\small \nopageexeed}
  \end{flushright}
\fi

%%%
%%% Dealing with page number excess
%%%
\def\npagemax{4}
\def\seclabelname{\ref{sec:citedreference}}
% Common to all the sections
\ifnum\value{page}>\npagemax%
  \vspace{1.0\baselineskip}
  \noindent
  \ifEnglishform
    \textcolor{red}{\textbf{\Large \fbox{Error!} Section {\seclabelname} exceeds the page limit (up to p.~{\npagemax}).}}
    \typeout{===ERROR===}
    \typeout{===ERROR=== Section \seclabelname exceeds the page limit (up to p.\npagemax).}
    \typeout{===ERROR===}
  \else
    \textcolor{red}{\textbf{\Large \fbox{警告} 第{\seclabelname}節がページ数の上限(p.~{\npagemax}まで)を超えています。(\thepage)}}
    \typeout{===警告===}
    \typeout{===警告=== 第\seclabelname 節がページ数の上限(p.\npagemax まで)を超えています。}
    \typeout{===警告===}
  \fi
\else
  % No problem in the page number limitation...
\fi

\clearpage
%%%%%%%%%%%%%%%%%%%%%%%%%%%%%%%%%%%%%%%%%%%%%%%%%%%%%%%%%%%%%%%%%%%%%%%
\ifEnglishform
  \section{Research results you have achieved so far\label{sec:pastwork}}
\else
  \section{これまでの研究成果\label{sec:pastwork}}
\fi

\iffinalform
% \input{4-achievement}
\else
  \ifEnglishform
    \textcolor{\remarkcolor}{If you have previous experience of using CfCA's open-use computational facilities, list the publications resulting from the activity.
\begin{itemize}
  \setlength{\itemsep}{-2pt}
  \item List only the publications of \underline{\textbf{your own research results.}}
  Do not bring references from other people's research that you cited in Sections \ref{sec:objective}, \ref{sec:method}, or \ref{sec:codeprep}.
  \item List separately the results you have obtained primarily through the use of CfCA's open-use computing facilities (Section \ref{pastusecfca}) and those obtained through the use of computing facilities at other institutions (Section \ref{pastuseother}).
  \item Provide as much detail as possible on the bibliographic information in the literature. Do not abbreviate the volume number or page numbers. % You may want to add DOI if you think you need.
  \item Include the specific names of the equipment and resources you used.
  \item The applicant's name should be written in bold in the list of authors.
  \item Use ``et al.'' or other abbreviations as appropriate when there are many authors.
\end{itemize}
Here are some examples.}
  \else
    \textcolor{\remarkcolor}{過去にCfCAの共同利用計算機や他機関の計算機設備等を利用した経験があれば、その成果による出版物を列挙してください。以下は注意事項です。
\begin{itemize}
  \setlength{\itemsep}{-2pt}
  \item ここに記すべきものは\underline{\textbf{申請者本人による研究成果}}です。第\ref{sec:objective}, \ref{sec:method}, \ref{sec:codeprep}節で引用された他研究者による先行研究の文献ではありません。
  \item CfCAの共同利用計算機を主に利用して得られた成果(第\ref{pastusecfca}節)と他機関の計算機設備等を利用して得られた成果(第\ref{pastuseother}節)を別個に記してください。
  \item 文献の書誌情報はなるべく詳しく記し、巻号ページ数も略さずに記載してください。 % DOIを付加して頂いても構いません。
  \item 利用した機材・資源の具体的な名称も書き入れてください。
  \item 著者一覧内で申請者本人の氏名は太字(ゴシック・ボールド体)で表記してください。
  \item 著者数が多い場合には``et al.''等を用いて適切に縮約してください。
\end{itemize}
以下は記入例です。}
  \fi % End of \ifEnglishform
\fi % End of \iffinalform

  \ifEnglishform
    \subsection{Publications achieved through the use of CfCA's facilities\label{pastusecfca}}
  \else
    \subsection{CfCAの共同利用計算機を利用した成果\label{pastusecfca}}
  \fi % End of \ifEnglishform
  \iffinalform
    \input{5-1-pastuse_cfca}
  \else
    \begin{itemize}
    \setlength{\itemsep}{-2pt}
    \item \textcolor{\remarkcolor}{\textbf{Iwasaki, K.}, Tomida, K., Inoue, T., Inutsuka, S. (2019) The early stage of molecular cloud formation by compression of two-phase atomic gases.\ \textit{The Astrophysical Journal,\/} volume 873, 6. (XC30, XC50)} % \textsf{10.3847/1538-4357/ab02ff}
    \item \textcolor{\remarkcolor}{Takasao, S., Tomida, K., \textbf{Iwasaki, K.}, Suzuki, T. K. (2019) Giant protostellar flares: Accretion-driven accumulation and reconnection-driven ejection of magnetic flux in protostars. \textit{The Astrophysical Journal,\/} volume 878, L10. (XC30, XC50)} % \textsf{10.3847/2041-8213/ab22bb}
    \item \textcolor{\remarkcolor}{Tsukamoto, Y., \textbf{Iwasaki, K.}, et al. (2015) Bimodality of circumstellar disk evolution induced by the hall current. \textit{The Astrophysical Journal,\/} volume 810, L26. (XC30)} % 10.1088/2041-8205/810/2/L26
    \end{itemize}
  \fi % End of \iffinalform

  \ifEnglishform
    \subsection{Publications achieved through the use of other facilities\label{pastuseother}}
  \else
    \subsection{他機関の計算機設備等を利用した成果\label{pastuseother}}
  \fi % End of \ifEnglishform
  \iffinalform
    \input{5-2-pastuse_other}
  \else
    \begin{itemize}
    \setlength{\itemsep}{-2pt}
    \item \textcolor{\remarkcolor}{\textbf{Iwasaki, K.} (2015) Minimizing dispersive errors in smoothed particle magnetohydrodynamics for strongly magnetized medium. \textit{Journal of Computational Physics,\/} volume 302, 359. (JCAHPC/Oakforest-PACS)} % \textsf{10.1016/j.jcp.2015.09.022}
    \end{itemize}
%
  \vfill
  \begin{flushright}
  \textcolor{\remarkcolor}{\small \nopageexeed}
  \end{flushright}
%
  \fi % End of \iffinalform

% \iffinalform
%   % Do nothing
% \else
%   \vfill
%   \begin{flushright}
%   \textcolor{\remarkcolor}{\small \nopageexeed}
%   \end{flushright}
% \fi % End of \iffinalform

%%%
%%% Dealing with page number excess
%%%
\def\npagemax{5}
\def\seclabelname{\ref{sec:pastwork}}
% Common to all the sections
\ifnum\value{page}>\npagemax%
  \vspace{1.0\baselineskip}
  \noindent
  \ifEnglishform
    \textcolor{red}{\textbf{\Large \fbox{Error!} Section {\seclabelname} exceeds the page limit (up to p.~{\npagemax}).}}
    \typeout{===ERROR===}
    \typeout{===ERROR=== Section \seclabelname exceeds the page limit (up to p.\npagemax).}
    \typeout{===ERROR===}
  \else
    \textcolor{red}{\textbf{\Large \fbox{警告} 第{\seclabelname}節がページ数の上限(p.~{\npagemax}まで)を超えています。(\thepage)}}
    \typeout{===警告===}
    \typeout{===警告=== 第\seclabelname 節がページ数の上限(p.\npagemax まで)を超えています。}
    \typeout{===警告===}
  \fi
\else
  % No problem in the page number limitation...
\fi

\clearpage
%%%%%%%%%%%%%%%%%%%%%%%%%%%%%%%%%%%%%%%%%%%%%%%%%%%%%%%%%%%%%%%%%%%%%%%
\ifXDS
  \setcounter{section}{5}  % Set it by hands!
  \ifEnglishform
    \section{The date when you wish to run your jobs\label{sec:preferredsched}}
  \else
    \section{ジョブの実行を希望する日程\label{sec:preferredsched}}
  \fi

  \iffinalform
   \input{6-schedule}
  \else
    \ifEnglishform
      \textcolor{\remarkcolor}{Indicate the dates on which you want to carry out the simulations if this application is adopted.}
    \else
      \textcolor{\remarkcolor}{本申請が採択された場合に、計算を実施する希望日程を記してください。}
    \fi
    \vfill
    \begin{flushright}
    \textcolor{\remarkcolor}{\small \nopageexeed}
    \end{flushright}
  \fi % End of iffinalform

%%%
%%% Dealing with page number excess
%%%
\def\npagemax{6}
\def\seclabelname{\ref{sec:preferredsched}}
% Common to all the sections
\ifnum\value{page}>\npagemax%
  \vspace{1.0\baselineskip}
  \noindent
  \ifEnglishform
    \textcolor{red}{\textbf{\Large \fbox{Error!} Section {\seclabelname} exceeds the page limit (up to p.~{\npagemax}).}}
    \typeout{===ERROR===}
    \typeout{===ERROR=== Section \seclabelname exceeds the page limit (up to p.\npagemax).}
    \typeout{===ERROR===}
  \else
    \textcolor{red}{\textbf{\Large \fbox{警告} 第{\seclabelname}節がページ数の上限(p.~{\npagemax}まで)を超えています。(\thepage)}}
    \typeout{===警告===}
    \typeout{===警告=== 第\seclabelname 節がページ数の上限(p.\npagemax まで)を超えています。}
    \typeout{===警告===}
  \fi
\else
  % No problem in the page number limitation...
\fi

\fi % End of ifXDS

%%%%%%%%%%%%%%%%%%%%%%%%%%%%%%%%%%%%%%%%%%%%%%%%%%%%%%%%%%%%%%%%%%%%%%%%%%%%
\end{document}
